\documentclass[11pt]{letter}
\usepackage[margin=1in]{geometry}
\usepackage[T1]{fontenc}
\usepackage{lmodern}
\usepackage{hyperref}

\signature{Daniel Kushnir and Hao Wang\\Rutgers University}
\address{Daniel Kushnir \quad \href{mailto:kushnir.daniel05@gmail.com}{kushnir.daniel05@gmail.com}\\
Hao Wang (corresponding author) \quad \href{mailto:hoguewang@gmail.com}{hoguewang@gmail.com}\\
Rutgers University}
\date{October 5, 2026}

\begin{document}
\begin{letter}{Editors-in-Chief\\Journal of Machine Learning Research}

\opening{Dear Editors,}

We submit our manuscript \emph{TSFM-Lens: Controlled, Causal and Confirmatory Interpretability for Time-Series
Foundation Models} for consideration as a regular paper in the Journal of Machine Learning Research.

Time-series foundation models are now used zero-shot for forecasting, but the interpretability methods developed for
language models do not carry over directly. Tokens span different lengths of time in different architectures, real
series carry no labels that could validate a probe, and the causal effects of internal features are small enough that
uncontrolled measurements produce confident false positives. The manuscript presents an open-source system built
around these three problems and the results of applying it to seven models.

The main contributions are:
\begin{enumerate}
  \item \textbf{A controlled setting.} A leakage-audited synthetic benchmark with exact generative ground truth,
  sealed dev and private splits, and a measured token-to-time alignment that puts architectures with tokens of 1 to
  96 time steps on a common axis.
  \item \textbf{Causal concept definitions with matched nulls.} A sparse-autoencoder feature counts as causal only if
  removing it changes the forecast more than removing the same per-token amount along a random direction on the same
  series, after reach controls confirm that the intervention lands. A planted-concept forecaster with known answers
  and decoys holds this test to a fixed false-positive gate; it showed our first null to be lenient.
  \item \textbf{One-shot confirmation.} Exploratory claims are frozen in a hash-pinned registry and tested once on the
  sealed split, with multiplicity handled per claim family. Of 82 replicable registered claims, 43 confirm.
  \item \textbf{Findings about seven models} (TimesFM, Chronos-2, Sundial, Chronos-Bolt, Chronos-T5, Timer and
  Time-MoE). Models agree on which inputs a concept groups (20 of 20 registered claims confirmed on 971 held-out
  series), but no cross-model claim about what a concept does to the forecast confirms, and no concept spans more than
  three of the seven models: models share what they select, not demonstrably what they do with it.
  \item \textbf{A record of corrected errors.} We document fifteen plausible results that the system's own controls
  showed to be wrong, as a guide for others working in this area.
\end{enumerate}

We believe the paper suits JMLR because its contribution is methodological as much as empirical: it argues, with
worked failures, that interpretability claims about forecasting models need known answers, matched nulls and held-out
confirmation, and it supplies a tool that makes these the default. All code, the benchmark generators, a complete
seven-model example run with its held-out confirmation, and a reproduction command for every reported number are
public under the MIT license at \url{https://github.com/DanK10097110/TSFM-Interp}.

The manuscript is not under review elsewhere. A preprint is available on arXiv at [arXiv identifier]. We plan to
describe the software itself separately in a short paper for the JMLR Machine Learning Open Source Software track; that
paper covers the toolkit's design and availability, not the scientific results reported here.

Use of AI tools: Claude (Anthropic) wrote most of the repository's code and tests, ran the experiments and helped draft
and edit the manuscript, under the authors' direction. The research questions, ideas and general design choices are
the authors'. The authors reviewed the implementation details Claude proposed, checked every reported result and take
full responsibility for the content.

Suggested action editors: [names]. Suggested reviewers: [names, affiliations, emails]. We have no conflicts of
interest to declare.

\closing{Sincerely,}

\end{letter}
\end{document}
